\subsection{WLAN-Anbindung eingebetteter Systeme} \todo{text}
% WLAN-Modus (Station Mode)
% IP-Adresse (DHCP)
% Verbindungsaufbau
% Netzwerkabhängigkeit
% lokale vs. Cloud-Systeme

% Die Einbindung des Systems in ein lokales Netzwerk erfolgt über die integrierte WLAN-Schnittstelle des Mikrocontrollers. Nach dem Verbindungsaufbau erhält das Gerät automatisch eine IP-Adresse und kann über diese im Netzwerk angesprochen werden. Dadurch wird eine direkte Kommunikation zwischen Nutzer und Gerät ohne zusätzliche Infrastruktur ermöglicht.